\section{Convergence Frameworks}\label{sec:convergence-frameworks}

This section shows how several mathematical frameworks relate to the same structural invariant once the optimizer quotient has been fixed. That invariant is structural rank. The bridge theorems below do not claim that all frameworks recover $\mathrm{srank}$ in the same formal sense: some give exact recovery, some give bounds, and some identify the same quotient branching regimes under exact decision preservation.

\subsection{Entropy and the Quotient}

The finite counting-to-probability bridge is established in Section~\ref{sec:counting-preliminaries} (Theorem~\ref{thm:bayes-from-counting}). Here we use that normalized layer to relate quotient entropy to structural rank.

\begin{theorem}[Entropy Bounded by Structural Rank]\label{thm:entropy-rank}
For Boolean coordinate spaces,
\[
H(\mathcal{D}) \le \mathrm{srank}(\mathcal{D}),
\]
where $H(\mathcal{D})$ is the base-$2$ entropy of the optimizer quotient.
\end{theorem}

\begin{proof}
Let $r=\mathrm{srank}(\mathcal{D})$. By definition, exactly $r$ coordinates are relevant, and this relevant set is sufficient (Section~\ref{sec:formal-setup}). Therefore $\Opt$ factors through the projection onto those $r$ coordinates:
\[
S \xrightarrow{\pi_R} \{0,1\}^r \xrightarrow{\bar{\Opt}} \mathcal{P}(A).
\]
Every optimizer class is therefore determined by one of the $2^r$ possible assignments to the relevant coordinates. Hence the optimizer quotient has at most $|\{0,1\}^r|=2^r$ classes. Under counting-normalized quotient entropy,
\[
H(\mathcal{D}) \le \log_2(2^r)=r=\mathrm{srank}(\mathcal{D}).
\]
\end{proof}

\subsection{Structural Rank Recovery}

\begin{theorem}[Fisher Support Recovers Structural Rank]\label{thm:fisher-rank-srank}
If one assigns score $1$ to relevant coordinates and $0$ to irrelevant ones, the resulting diagonal Fisher-style support matrix has rank $\mathrm{srank}(\mathcal{D})$.
\end{theorem}

\begin{proof}
In the stated encoding, the diagonal entry for coordinate $i$ is nonzero exactly when $i$ is relevant. For any diagonal matrix, rank equals the number of nonzero diagonal entries. Therefore the rank is exactly
\[
\big|\{i:\ i \text{ relevant}\}\big|=\mathrm{srank}(\mathcal{D}).
\]
\end{proof}

\begin{theorem}[Zero-Distortion Summaries Are Constrained by Quotient Structure]\label{thm:rate-distortion-bridge}
At zero distortion, any lossless summary that preserves optimal actions must assign distinct codewords to distinct optimizer-quotient classes. Thus the admissible summary size is constrained by the same quotient structure determined by the relevant-coordinate support.
\end{theorem}

\begin{proof}
At zero distortion, any two states mapped to the same summary must induce the same optimal-action set; otherwise optimal actions are not preserved exactly. So admissible summaries are precisely coarsenings that refine decision equivalence, i.e., they cannot identify states across distinct optimizer-quotient classes.

Therefore any lossless summary alphabet must contain at least one distinct codeword for each optimizer-quotient class. If the alphabet were smaller than the number of decision classes, two different classes would share a codeword, violating zero distortion. Conversely, assigning distinct codewords to the classes yields an exact lossless summary.

By Section~\ref{sec:formal-setup}, the quotient class structure is determined by the relevant-coordinate support, hence by $\mathrm{srank}(\mathcal{D})$. So the zero-distortion summary size is constrained by the same support-size core.
\end{proof}

\begin{theorem}[Wasserstein Regimes Follow Quotient Branching]\label{thm:wasserstein-bridge}
When transport is computed on the optimizer-quotient classes, transport cost is zero in the singleton-class regime and strictly positive once multiple optimizer classes must be distinguished. Thus transport complexity tracks the same quotient branching regimes determined by decision relevance.
\leanmeta{\LHrng{W}{1}{4}.}
\end{theorem}

\begin{proof}
If the quotient has one class, the identity coupling is diagonal and incurs zero transport cost. If the quotient has at least two classes with nontrivial mass split, any coupling that matches the distinct targets carries off-diagonal mass and therefore positive transport cost. By Section~\ref{sec:formal-setup}, the same class multiplicity is controlled by relevant-coordinate support, so the transport regimes are determined by the same structural core.
\end{proof}

\begin{theorem}[Decision-Relevant Information Requires Nontriviality]\label{thm:nontriviality-counting}
If a decision problem carries nonzero decision-relevant information, then its optimizer quotient is nontrivial. Equivalently, a singleton quotient cannot support positive decision-relevant entropy.
\end{theorem}

\begin{proof}
If the optimizer quotient is singleton, then all states are decision-equivalent and $\Opt$ is constant on $S$. Hence there is only one attained optimizer value and the quotient entropy is
\[
H(\mathcal{D})=\log_2(1)=0.
\]
So positive decision-relevant information is impossible in the singleton case. Taking contrapositive yields the theorem.
\end{proof}

\subsection{Framework Convergence}

\begin{theorem}[Six-Framework Comparison Around Structural Rank]\label{thm:six-framework-recovery}
Within the finite setting of this paper, the following objects are linked by the optimizer quotient and the relevant-coordinate support:
\begin{enumerate}
\item relevant-coordinate support,
\item structural rank,
\item optimizer-quotient image size,
\item quotient entropy,
\item Fisher-style support count,
\item zero-distortion lossless decision summary size.
\end{enumerate}
\leanmeta{\LH{SK1}, \LH{QT5}, \LH{IT3}, \LH{FS2}, \LH{RS5}.}
\end{theorem}

\begin{proof}
The chain is explicit:
\begin{enumerate}
\item Section~\ref{sec:formal-setup} identifies the support core: relevant coordinates and structural rank (Definition~\ref{def:srank}; Proposition~\ref{prop:minimal-relevant-equiv}).
\item Proposition~\ref{prop:optimizer-coimage} and Theorem~\ref{thm:quotient-universal} identify optimizer-quotient classes as the canonical decision abstraction and relate them to the image of $\Opt$.
\item Theorem~\ref{thm:entropy-rank} bounds quotient entropy in terms of this support size.
\item Theorem~\ref{thm:fisher-rank-srank} shows Fisher-style support counting recovers exactly the same cardinality.
\item Theorem~\ref{thm:rate-distortion-bridge} shows zero-distortion summary size is constrained by the same quotient class structure.
\item Theorem~\ref{thm:nontriviality-counting} excludes degenerate singleton quotients when decision-relevant information is positive.
\end{enumerate}
Thus the six listed objects are organized around the same support/quotient core, with exact recovery in some cases and structural bounds in others.
\end{proof}

The point is not that these frameworks become identical in every technical sense. Rather, once the optimizer quotient is fixed, they can all be compared against the same support/quotient core for decision relevance. The paper's backbone is therefore simple: shared definitions, the quotient universal property, and a family of bridge theorems centered on $\mathrm{srank}$.
